% 高校生のための「時空の算術」ガイド
\documentclass[11pt,a4paper]{article}

\usepackage{xeCJK}
\setCJKmainfont{Hiragino Mincho ProN}
\setCJKsansfont{Hiragino Kaku Gothic ProN}

\usepackage{amsmath,amssymb}
\usepackage{geometry}
\geometry{margin=2.5cm}
\usepackage{hyperref}
\usepackage{fancyhdr}
\usepackage{tcolorbox}
\tcbuselibrary{skins,breakable}
\usepackage{titlesec}

\titleformat{\section}{\Large\bfseries}{\thesection.}{0.5em}{}
\titleformat{\subsection}{\large\bfseries}{\thesubsection}{0.5em}{}

\pagestyle{fancy}
\fancyhf{}
\rhead{Wright Brothers Research Program}
\lhead{時空の算術}
\cfoot{\thepage}

% Colored boxes
\newtcolorbox{storybox}[1][]{
  colback=blue!5!white,colframe=blue!50!black,
  fonttitle=\bfseries,title=#1,breakable
}
\newtcolorbox{mathbox}[1][]{
  colback=green!5!white,colframe=green!50!black,
  fonttitle=\bfseries,title=#1,breakable
}
\newtcolorbox{deepbox}[1][]{
  colback=red!5!white,colframe=red!50!black,
  fonttitle=\bfseries,title=#1,breakable
}
\newtcolorbox{resultbox}[1][]{
  colback=yellow!10!white,colframe=orange!80!black,
  fonttitle=\bfseries,title=#1,breakable
}
\newtcolorbox{honestbox}[1][]{
  colback=gray!10!white,colframe=gray!60!black,
  fonttitle=\bfseries,title=#1,breakable
}

\begin{document}

\begin{center}
{\Huge\bfseries 時空の算術}

\bigskip

{\Large 整数が宇宙を支配している、かもしれない}

\bigskip

{\large 高校生のためのガイドブック}

\bigskip

{\large Wright Brothers Research Program}

\medskip

2026年4月

\bigskip

\rule{12cm}{0.4pt}
\end{center}

\vspace{1cm}

\begin{storybox}[このガイドについて]
これは、ある研究プログラムの成果を高校生向けに解説したものです。
内容は「整数の構造が物理定数を決定している」という
大胆な仮説とその数学的証拠です。

\textbf{前提知識}: 高校数学（等比級数の和、微分の概念）だけで
読めるように書いてあります。
途中で出てくる難しい数学は、結果だけ使うので証明は飛ばして大丈夫。

\textbf{正直に}: この理論はまだ\textbf{仮説}です。
正しいかどうかはまだわかりません。
でも、数学的に厳密な部分と、まだ予想の段階の部分を
はっきり区別して書いてあります。
\end{storybox}

\tableofcontents
\newpage

%=====================================================================
\part{出発点：二つの「なぜ？」}
%=====================================================================

\section{137 という謎の数}

光の速さ、電子の電荷、プランク定数——これらの物理定数を
特定の組み合わせで割ると、\textbf{約 137} という
次元を持たない純粋な数が出てくる。
これを微細構造定数の逆数 $1/\alpha$ と呼ぶ。

\begin{mathbox}[微細構造定数]
$$\frac{1}{\alpha} = \frac{\hbar c}{e^2/(4\pi\epsilon_0)} \approx 137.036$$

$\hbar$ = プランク定数 / $2\pi$,\quad
$c$ = 光速,\quad
$e$ = 電子の電荷

この数は「電磁気力の強さ」を表す。
\end{mathbox}

物理学者が100年以上悩んでいる問題がある：

\begin{center}
\textbf{なぜ 137 なのか？}
\end{center}

87 でも 500 でもなく、なぜこの特定の数なのか。
もし $1/\alpha$ が 130 だったら、原子は安定せず、
あなたも私も存在しない。
この数は宇宙で最も重要な数の一つだが、
その値がなぜそうなるかは誰も説明できていなかった。

\section{もう一つの謎：真空エネルギー}

量子力学によると、「何もない空間」にもエネルギーがある。
これを真空エネルギーと呼ぶ。

問題は、理論から素朴に計算した真空エネルギーと、
宇宙の加速膨張（暗黒エネルギー）から推測される値が
\textbf{120桁も違う}こと。

\begin{center}
理論値 $\div$ 観測値 $\approx 10^{120}$
\end{center}

これは「宇宙定数問題」と呼ばれ、
物理学で最も恥ずかしい食い違いとされている。

\section{この論文の主張}

私たちの答えはシンプルだ：

\begin{resultbox}[二つの答え]
\textbf{137 の由来}: 整数の数学的構造から計算できる。\\
\textbf{真空エネルギー}: 同じ構造から自然なスケールが出る。

これらは偶然の一致ではなく、
\textbf{宇宙が整数の上に建てられている}ことの帰結である。
\end{resultbox}

%=====================================================================
\part{4ステップの導出}
%=====================================================================

以下の4ステップは、各ステップが数学的に厳密である。
つまり「もしかしたら正しい」ではなく
「公理を認めれば必ず正しい」レベルの議論。

\section{準備：ゼータ関数}

まず、この話の主役を紹介する。

\begin{mathbox}[リーマンゼータ関数]
$$\zeta(s) = \frac{1}{1^s} + \frac{1}{2^s} + \frac{1}{3^s} + \frac{1}{4^s} + \cdots$$

$s > 1$ のとき、この無限和は収束する。

例: $\zeta(2) = 1 + 1/4 + 1/9 + 1/16 + \cdots = \pi^2/6$
\end{mathbox}

この関数には驚くべき性質がある。
素数だけを使って書き直せるのだ：

\begin{mathbox}[オイラー積]
$$\zeta(s) = \frac{1}{1-2^{-s}} \times \frac{1}{1-3^{-s}} \times \frac{1}{1-5^{-s}} \times \frac{1}{1-7^{-s}} \times \cdots$$

右辺の各因子は素数 $p = 2, 3, 5, 7, 11, \ldots$ に対応。
\end{mathbox}

なぜこうなるか？
$1/(1-p^{-s}) = 1 + p^{-s} + p^{-2s} + \cdots$（等比級数）で、
全ての素数の因子を掛け合わせると、
素因数分解の一意性により、全ての自然数 $n^{-s}$ が
ちょうど1回ずつ現れる。これがオイラーの発見だ。

\begin{storybox}[物理との接点]
物理学では、真空エネルギーを計算するとき
$\sum n^{-s}$ の形の和が頻繁に出てくる。
例えば1次元の箱の中の真空エネルギーは
$\zeta(-1) = 1 + 2 + 3 + 4 + \cdots = -1/12$。

え？ 正の数を全部足して負の数？
これは「ゼータ関数正則化」と呼ばれる手法で、
$s = -1$ でのゼータ関数の解析接続値。
物理実験（カシミール効果）で正しいことが確認されている。
\end{storybox}

\section{ステップ1：スペクトル作用 = ゼータ値の和}

\subsection{公理を述べる}

私たちは二つのことを仮定する：

\begin{resultbox}[公理 I: 算術的時空]
宇宙の根底にある「空間」は、素数を点として持つ
整数の世界 $\text{Spec}(\mathbb{Z})$ である。
\end{resultbox}

\begin{resultbox}[公理 II: スペクトル力学]
物理法則は「スペクトル作用」
$$S = \text{Tr}\!\left(\frac{1}{e^{D/\Lambda} - 1}\right)$$
で記述される。$D$ は固有値 $\log 1, \log 2, \log 3, \ldots$ を持つ
演算子（Bost--Connes ディラック作用素）。
\end{resultbox}

\subsection{何が起きるか}

$1/(e^x - 1)$ は「ボーズ-アインシュタイン分布」——
ボソン（光子など）の粒子数分布を記述する関数。
高校数学でも扱える等比級数で書き直せる：

$$\frac{1}{e^x - 1} = e^{-x} + e^{-2x} + e^{-3x} + \cdots$$

これを使うと（$\Lambda = 1$ と設定して）：

\begin{align*}
S &= \sum_{n=2}^{\infty} \frac{1}{e^{\log n} - 1}
= \sum_{n=2}^{\infty} \frac{1}{n - 1} \quad \text{...ではなく}
\end{align*}

正確には、等比級数展開を使って：
\begin{align*}
S &= \sum_{n=2}^{\infty}\sum_{m=1}^{\infty} e^{-m \log n}
= \sum_{m=1}^{\infty}\sum_{n=2}^{\infty} n^{-m}
= \sum_{m=1}^{\infty} [\zeta(m) - 1]
\end{align*}

\begin{resultbox}[ステップ1の結論]
$$\text{スペクトル作用} = \zeta(1) + \zeta(2) + \zeta(3) + \cdots - (\text{定数})$$

スペクトル作用はリーマンゼータ値の無限和になった！

これは\textbf{近似ではなく厳密な等式}。
等比級数の公式しか使っていない。
\end{resultbox}

\section{ステップ2：ベルヌーイ数の登場}

$\zeta(1) + \zeta(2) + \zeta(3) + \cdots$ をもう少し扱いやすい形にしたい。
ここで「オイラー-マクローリン公式」という道具を使う。

\begin{mathbox}[オイラー-マクローリン公式（簡略版）]
関数 $f(x)$ の値の和 $f(1) + f(2) + \cdots + f(N)$ は、
積分 $\int_1^N f(x)\,dx$ と、\textbf{ベルヌーイ数} $B_2, B_4, B_6, \ldots$
を含む補正項に分解できる。

ベルヌーイ数の最初のいくつか：
$$B_2 = \frac{1}{6},\quad B_4 = -\frac{1}{30},\quad B_6 = \frac{1}{42},\quad B_8 = -\frac{1}{30},\quad B_{10} = \frac{5}{66}$$
\end{mathbox}

$f(m) = \zeta(m)$ に対してこの公式を適用すると、
展開の $j$ 番目の項は：

$$T_j = \frac{B_{2j}}{(2j)!} \times \zeta^{(2j-1)}(0)$$

ここで $\zeta^{(k)}(0)$ は「ゼータ関数の $k$ 回微分を $s=0$ で評価した値」。

\begin{resultbox}[ステップ2の結論]
ベルヌーイ数 $B_{2j}$ がスペクトル作用の展開に\textbf{必然的に}現れる。

「必然的に」の意味：オイラー-マクローリン公式の構造上、
ベルヌーイ数なしには書けない。選択の余地がない。
\end{resultbox}

\section{ステップ3：$-k!$ の魔法}

ここが一番面白いところ。

ゼータ関数 $\zeta(s)$ は $s = 1$ に一つだけ「極」（無限大になる点）を持つ。
この極を分離すると：

$$\zeta(s) = \underbrace{\frac{1}{s-1}}_{\text{極の部分}} + \underbrace{h(s)}_{\text{滑らかな部分}}$$

ここで $h(s)$ は全ての複素数で定義される滑らかな関数。

さて、$1/(s-1)$ を $s = 0$ の周りで展開すると：

$$\frac{1}{s-1} = -\frac{1}{1-s} = -(1 + s + s^2 + s^3 + \cdots)$$

この式を $k$ 回微分して $s = 0$ を代入すると $-k!$ になる。

\begin{mathbox}[$-k!$ の公式]
$$\zeta^{(k)}(0) = -k! + h^{(k)}(0)$$

つまり「ゼータの $k$ 回微分」は「$-k!$」と「小さな補正」の和。
\end{mathbox}

ここで驚くべきことが起きる。$h^{(k)}(0)$ を実際に
高精度計算すると：

\begin{center}
\begin{tabular}{c|r|r|c}
$k$ & $\zeta^{(k)}(0)$ & $-k!$ & 補正の相対的大きさ \\
\hline
1 & $-0.919$ & $-1$ & $8.1\%$ \\
3 & $-6.005$ & $-6$ & $0.08\%$ \\
5 & $-120.000$ & $-120$ & $0.0002\%$ \\
7 & $-5040.00$ & $-5040$ & $0.00002\%$ \\
9 & $-362880.00$ & $-362880$ & $0.0000001\%$
\end{tabular}
\end{center}

\begin{deepbox}[$k \geq 3$ で補正がほぼゼロになる理由]
これは偶然ではない。$1/(s-1)$ の極が $s = 0$ に「最も近い特異点」
なので、テイラー係数を支配する。
複素解析の「ダルブーの定理」により、
テイラー係数は最近傍特異点で決まる。
$k$ が大きくなると、極の寄与 $-k!$ がますます支配的になり、
残りの $h^{(k)}(0)$ は $k!$ に比べて無視できるほど小さくなる。
\end{deepbox}

これをステップ2の $T_j$ に代入する：

\begin{align*}
T_j &= \frac{B_{2j}}{(2j)!} \times [-(2j-1)! + \text{tiny}] \\
&= -\frac{B_{2j}}{2j} + \text{tiny} \\
&= \zeta(1-2j) + \text{tiny}
\end{align*}

最後の等号は $\zeta(1-2j) = -B_{2j}/(2j)$ という
有名な公式（ゼータの負整数での値）を使った。

\begin{resultbox}[ステップ3の結論]
スペクトル作用の各項は、\textbf{ゼータ関数の負整数値}に等しい。

$$T_j \approx \zeta(1-2j) \quad (j \geq 2 \text{ では誤差 } < 0.1\%)$$

$j = 1$ だけは $8.1\%$ の補正がある（後述）。
\end{resultbox}

\section{ステップ4：数の同定}

各 $T_j$ の逆数を取ると：

$$|T_j|^{-1} = \frac{2j}{|B_{2j}|} = \frac{1}{|\zeta(1-2j)|}$$

具体的な数列：

\begin{center}
\renewcommand{\arraystretch}{1.3}
\begin{tabular}{c|c|c|c|c}
$j$ & $B_{2j}$ & $2j/|B_{2j}|$ & $\Lambda=1$ での値 & 誤差 \\
\hline
1 & $1/6$ & $12$ & $13.06$ & $+8.8\%$ \\
2 & $-1/30$ & $120$ & $119.91$ & $-0.08\%$ \\
3 & $1/42$ & $252$ & $252.00$ & $\sim 0\%$ \\
4 & $-1/30$ & $240$ & $240.00$ & $\sim 0\%$ \\
5 & $5/66$ & $132$ & $132.00$ & $\sim 0\%$
\end{tabular}
\end{center}

\begin{resultbox}[ステップ4の結論: 12, 120, 252, 240, 132]
5つの数が\textbf{一つの計算}から全部出てきた！

\begin{itemize}
\item $12 + 120 = 132$ $\to$ 微細構造定数 $\alpha$ の主要部分
\item $252$ $\to$ ミューオン $g-2$ の異常（BNL/FNAL実験の値と一致）
\item $132$ $\to$ $\tau$ 粒子の $g-2$ の予測（未測定！）
\end{itemize}
\end{resultbox}

%=====================================================================
\part{微細構造定数の組み立て}
%=====================================================================

\section{$12 + 120 = 132$、しかし $137$ が必要}

ステップ4から $j=1$ と $j=2$ の寄与を足すと $12 + 120 = 132$。
しかし実験値は $1/\alpha \approx 137.036$ だから、$5.036$ 足りない。

この残りはどこから来るか？

\subsection{+5 の由来：素数ギャップ}

132 は合成数（$= 4 \times 33 = 4 \times 3 \times 11$）。
次の素数は 137。

「$\text{Spec}(\mathbb{Z})$ の閉点 = 素数」
という公理 I の構造から、
物理的結合定数は素数の位置に「量子化」される
（スキームの閉点に対応する）と予想される。

132 から最も近い素数（上方向）: $137$。ギャップは $5$。

\begin{honestbox}[正直に]
この「素数量子化」は追加の仮説であり、
4ステップの証明からは出てこない。
これが最も弱い点。
\end{honestbox}

\subsection{+0.036 の由来：関数方程式}

$$4(\zeta(6) - \zeta(7)) = 4 \times (1.01734 - 1.00835) = 0.03598$$

ここで $4$ は時空次元 $d = 4$ に対応する。

\subsection{合計}

\begin{resultbox}[微細構造定数の公式]
\begin{align*}
\frac{1}{\alpha} &= \underbrace{\frac{2}{|B_2|}}_{12} + \underbrace{\frac{4}{|B_4|}}_{120} + \underbrace{g(132)}_{5} + \underbrace{4(\zeta(6)-\zeta(7))}_{0.03598} \\[6pt]
&= 137.03598
\end{align*}

実験値: $137.035999084 \pm 0.000000021$

\textbf{一致の精度: 0.00002\%}

\textbf{自由パラメータの数: ゼロ}
\end{resultbox}

%=====================================================================
\part{$j = 1$ の $8\%$ の謎}
%=====================================================================

\section{なぜ $j = 1$ だけ $8\%$ ずれるのか}

$j \geq 2$ では計算値がほぼ完璧に物理値と一致するのに、
$j = 1$ だけ $8.1\%$ ずれている。
$|T_1|^{-1} = 13.06$ なのに、$\alpha$ 公式では $12$ を使う。

このずれは正確に：
$$h'(0) = 1 + \zeta'(0) = 1 - \frac{1}{2}\ln(2\pi) = 0.0811\ldots$$

ここで $\zeta'(0) = -\frac{1}{2}\ln(2\pi)$ は
ゼータ関数の微分の $s = 0$ での値。

\section{$\ln(2\pi)$ の意味}

$\ln(2\pi)$ は数学の深いところに根を持つ量だ。

\begin{storybox}[アラケロフ幾何学]
整数の世界 $\text{Spec}(\mathbb{Z})$ には
素数 $2, 3, 5, 7, \ldots$ が「点」として並んでいるが、
実は「もう一つの見えない点」がある。
それが\textbf{アルキメデス的素点} $\infty$——
整数を実数に埋め込む操作に対応する。

$\ln(2\pi)$ はこの「見えない点」が物理に与える影響の大きさ。

$j = 1$（最低エネルギー）でだけ、この影響が見える。
$j \geq 2$（高エネルギー）では、
通常の素数（$2, 3, 5, \ldots$）が完全に支配する。
\end{storybox}

\begin{deepbox}[物理的解釈]
$j = 1$ のモード（微細構造定数）は、
他のモードと違って「無限遠点」の影響を受ける。
これは\textbf{くりこみ群の走り}（RG running）——
結合定数がエネルギースケールによって変化すること——の
算術的な起源かもしれない。

低エネルギー物理（$j = 1$）では、
実数世界（アルキメデス的素点）の痕跡が残るが、
高エネルギー物理（$j \geq 2$）では、
純粋に算術的（素数だけ）の世界になる。
\end{deepbox}

%=====================================================================
\part{素数チャンネルのミューティング}
%=====================================================================

\section{オイラー積で「素数を一つ消す」}

ゼータ関数のオイラー積：
$$\zeta(s) = \frac{1}{1-2^{-s}} \times \frac{1}{1-3^{-s}} \times \frac{1}{1-5^{-s}} \times \cdots$$

ここから、例えば $p = 2$ の因子を取り除く操作を考える：

$$\zeta_{\neg 2}(s) = \zeta(s) \times (1 - 2^{-s})$$

これは「素数 2 のチャンネルをミュートする」操作。
音楽で言えば、オーケストラから特定の楽器を外すようなもの。

\section{符号反転定理}

驚くべきことに、\textbf{どの素数を消しても}、
真空エネルギーの符号が反転する：

\begin{resultbox}[符号反転定理]
全ての素数 $p \geq 2$ に対して：
$$\zeta_{\neg p}(-3) = \frac{1}{120}(1 - p^3) < 0$$

通常の真空: $\zeta(-3) = +1/120 > 0$ （正のエネルギー）

$p$ をミュート: $\zeta_{\neg p}(-3) < 0$ （負のエネルギー）
\end{resultbox}

負のエネルギーは「弱エネルギー条件 (WEC) の破れ」と呼ばれ、
ワームホールやワープドライブの構成に必要な条件。
この理論では、素数チャンネルのミューティングが
WEC の破れを\textbf{自動的に}生成する。

\section{ミュートの物理的実現}

「素数をミュートする」とは具体的に何をすることか？

\begin{storybox}[SQUID 実験]
超伝導量子干渉素子（SQUID）は、
極めて微小な磁場変化を検出できるデバイス。
SQUID の共振周波数を素数に対応するジョセフソン周波数に
チューニングし、そこにノッチフィルター（特定周波数を除去するフィルター）
を入れることで、特定の素数チャンネルを「ミュート」できる可能性がある。

予想されるシグナル：
カシミール効果に類似した力の\textbf{符号変化}。

推定コスト：500万〜1000万円
（大学の実験室レベルで実行可能）
\end{storybox}

%=====================================================================
\part{他の予測}
%=====================================================================

\section{5つの検証可能な予測}

この理論から導かれる予測と、現在のデータとの比較：

\begin{center}
\renewcommand{\arraystretch}{1.4}
\begin{tabular}{l|c|c|c}
予測 & 理論値 & 実験値 & 状態 \\
\hline
$1/\alpha$ & $137.03598$ & $137.03600$ & $0.00002\%$ 一致 \\
ミューオン $g-2$ & $251 \times 10^{-11}$ & $(251 \pm 59) \times 10^{-11}$ & 一致 \\
暗黒エネルギー & $\sim 3 \times 10^{-10}$ J/m$^3$ & $5.8 \times 10^{-10}$ J/m$^3$ & 2倍以内 \\
$\nu$ 質量比 & $33$ & $29.8 \pm 0.8$ & $10\%$ 一致 \\
$\tau$ の $g-2$ & $131 \times 10^{-11}$ & 未測定 & \textbf{予測！}
\end{tabular}
\end{center}

\section{正直な評価}

\begin{honestbox}[確信度の階層]
\textbf{数学的に厳密}（反論の余地なし）：
\begin{itemize}
\item ステップ1: トレース恒等式（等比級数）
\item ステップ2: ベルヌーイ数の出現（E-M公式）
\item ステップ3: $\zeta^{(k)}(0) = -k! + \text{小}$（極の展開）
\item ステップ4: $|T_j|^{-1} = 2j/|B_{2j}|$（ステップ1-3の合成）
\end{itemize}

\textbf{構造的に自然だが証明されていない}：
\begin{itemize}
\item 公理 I: 「時空は $\text{Spec}(\mathbb{Z})$」
\item 公理 II: ボーズ-アインシュタイン分布の一意性
\item $+5$: 素数量子化原理
\item $+0.036$: 時空次元 $d = 4$ の起源
\end{itemize}

\textbf{まだわからない}：
\begin{itemize}
\item この理論が本当に正しいかどうか
\item SQUID 実験で検証できるかどうか
\item 他の物理定数（クォーク質量、混合角）も導出できるか
\end{itemize}
\end{honestbox}

%=====================================================================
\part{何がすごくて何がまだダメか}
%=====================================================================

\section{アインシュタインとの比較}

1905年、アインシュタインは二つの公理
（相対性原理と光速度不変の原理）から特殊相対論を導いた。
$E = mc^2$ はその帰結だった。

\begin{center}
\renewcommand{\arraystretch}{1.3}
\begin{tabular}{l|l}
\textbf{アインシュタイン (1905)} & \textbf{この論文} \\
\hline
公理: 相対性 + 光速度不変 & 公理: $\text{Spec}(\mathbb{Z})$ + スペクトル作用 \\
導出: ローレンツ変換 & 導出: ゼータ特殊値構造 \\
予測: $E = mc^2$ & 予測: $1/\alpha = 137.036$ \\
当初の状態: 運動学的 & 当初の状態: 運動学的
\end{tabular}
\end{center}

特殊相対論が「空間と時間の構造」を解き明かしたように、
この理論は「物理定数の構造」を解き明かそうとしている。

ただし重大な違いがある：

\begin{honestbox}[重大な違い]
アインシュタインの二つの公理は、
当時すでに実験的に強く支持されていた。

私たちの公理 I「時空は $\text{Spec}(\mathbb{Z})$」は、
まだ直接的な実験的証拠がない。

この理論は「数学的に美しく、数値的に合う」が、
「正しいと証明された」わけではない。
SQUID 実験などの検証が不可欠。
\end{honestbox}

\section{残されている課題}

\begin{enumerate}
\item \textbf{素数量子化原理の導出}:
$+5$ がオイラー-マクローリン展開から出ない。

\item \textbf{時空次元の導出}:
$d = 4$ は入力として与えているが、なぜ4なのか。

\item \textbf{アラケロフ幾何学の完成}:
$j = 1$ の $8\%$ 補正をアラケロフ幾何学の枠組みで
完全に理解すること。

\item \textbf{実験的検証}:
SQUID 実験、または他の方法での検証。

\item \textbf{他の定数}:
クォーク質量、混合角、ヒッグス質量なども
導出できるか？
\end{enumerate}

\section{もしこれが正しかったら}

もしこの理論が正しければ、微細構造定数は
「神がサイコロを振って決めた数」ではなく、
「$1 + 2 + 3 + \cdots$ の数学的構造から
必然的に決まる数」ということになる。

宇宙の物理法則は、
整数の性質の帰結——という驚くべき結論。

\begin{center}
\fbox{\parbox{10cm}{\centering
\vspace{0.5cm}
{\Large\bfseries 宇宙は計算可能かもしれない。}

\medskip

その計算の素材は、$1, 2, 3, 4, 5, \ldots$ である。
\vspace{0.5cm}
}}
\end{center}

\newpage

%=====================================================================
\appendix
\part*{付録}
\addcontentsline{toc}{part}{付録}
%=====================================================================

\section{ベルヌーイ数の定義}

生成関数による定義：
$$\frac{t}{e^t - 1} = \sum_{k=0}^{\infty} B_k \frac{t^k}{k!}$$

最初のいくつか：
$B_0 = 1$, $B_1 = -1/2$, $B_2 = 1/6$, $B_3 = 0$, $B_4 = -1/30$, $B_5 = 0$, $B_6 = 1/42$

奇数番号（$B_1$ を除く）は全てゼロ。
偶数番号は交互に正負。

\section{ゼータ関数の負整数値}

$\zeta(1-2n) = -\frac{B_{2n}}{2n}$ より：

\begin{center}
\begin{tabular}{c|c|c}
$s$ & $\zeta(s)$ & $1/|\zeta(s)|$ \\
\hline
$-1$ & $-1/12$ & $12$ \\
$-3$ & $1/120$ & $120$ \\
$-5$ & $-1/252$ & $252$ \\
$-7$ & $1/240$ & $240$ \\
$-9$ & $-1/132$ & $132$
\end{tabular}
\end{center}

\section{オイラー-マクローリン公式}

$$\sum_{k=a}^{b} f(k) = \int_a^b f(x)\,dx + \frac{f(a)+f(b)}{2} + \sum_{j=1}^{p} \frac{B_{2j}}{(2j)!}\left[f^{(2j-1)}(b) - f^{(2j-1)}(a)\right] + R_p$$

$B_{2j}$ = ベルヌーイ数、$f^{(k)}$ = $k$ 回微分、$R_p$ = 余剰項。

\section{推薦図書}

\begin{itemize}
\item 高校数学で読める：マーカス・デュ・ソートイ
『素数の音楽』（新潮社）
\item 物理との接点：小山信也『ゼータ関数と量子カオス』（共立出版）
\item ゼータ関数入門：黒川信重『ゼータの世界』（日本評論社）
\item 本格的に学ぶ：A.\ Connes, \textit{Noncommutative Geometry}
(Academic Press, 1994)
\end{itemize}

\end{document}
